\documentclass[10pt,a4paper]{amsart}
\usepackage[margin=19mm]{geometry}
\usepackage{amsmath,amssymb,mathtools}
\usepackage[scr=rsfs,cal=euler]{mathalfa}
\usepackage{xfrac}
\usepackage[unicode,hidelinks,bookmarks=false]{hyperref}
\newcommand{\ie}{i.e.\ }
\newcommand{\cf}{cf.\ }
\newcommand{\resp}{respectively\ }
\newcommand{\Subsection}{\subsection}
\providecommand{\nobreakdash}{\nobreak\hbox{-}}
\setlength{\emergencystretch}{2em}
\multlinegap0pt
\usepackage{bm}
\usepackage{bbm}
\usepackage{dsfont}
\usepackage{xspace}
\usepackage{cancel}
\usepackage{mathtools}
\usepackage{amsthm}
\makeatletter
\@ifundefined{theorem}{\newtheorem{theorem}{Theorem}}{}
\@ifundefined{lemma}{\newtheorem{lemma}[theorem]{Lemma}}{}
\makeatother
\newcommand{\ISO}{\mathrm{ISO}}
\newcommand{\One}[1]{{\mathbbm{1}}\left\{{#1}\right\}}
\newcommand{\R}{\mathbb{R}}
\newcommand{\bZ}{\mathbf{Z}}
\title[Testing conditional independence under isotonicity]{Testing conditional independence under isotonicity}
\author{Rohan Hore, Jake A. Soloff, Rina Foygel Barber, Richard J. Samworth}
\date{Source: arXiv:2501.06133v3 (CC BY 4.0)}
\begin{document}
\maketitle

\noindent\emph{Source and licence.} Testing conditional independence under isotonicity. Version arXiv:2501.06133v3 (CC BY 4.0), distributed under the Creative Commons Attribution 4.0 International licence, as recorded in the arXiv licence metadata for this exact version and shown on the abstract page https://arxiv.org/abs/2501.06133v3. Full rights notice: licenses/arxiv-2501.06133-CC-BY-4.0.txt.\par\smallskip

\noindent\textit{Editorial scope.} Counted unit: the Lemma whose source label is \texttt{lem:characterizing-isotonic-L1-projection\_} in the article (source0.tex lines 1859--1868, source label \texttt{lem:characterizing-isotonic-L1-projection\_}), delivered together with the article's own complete original proof of it (source0.tex lines 1869--1970, one \texttt{proof} environment of 102 lines). No mathematics is written, repaired or re-derived: statement and proof are the article's own text line for line. Required context retained uncounted: Eq:utildeiso (lines 1853--1855); lem:a-median-result (lines 2510--2512). Every definition, display and label of those lines is retained unchanged. Omitted from this package and not redistributed: 1--1852, 1856--1858, 1971--2509, 2513--2838. Nothing inside a retained excerpt is omitted or altered: every retained body is delivered exactly as the article's lines are written, so no omission receipt of any kind accompanies this package. Citations: the retained excerpts carry no citation command at all, so no bibliography entry of the article is transcribed and no reference is redistributed. Reference adaptation: none was needed and none was made; every reference inside the delivered unit resolves inside it, so no unresolved reference or citation is produced. Printed slips kept literal: no printed word, symbol, hypothesis or number of the counted statement or of its proof was altered. Layout and class: the article's own class is \texttt{article} with packages including amsmath, amssymb, amsthm, bbm, natbib, xcolor, tikz, algorithm, algorithmic, mathtools, calc, enumerate, parskip, longtable; the wrapper is the lane's pinned \texttt{amsart} editorial wrapper at 10pt on A4, which restates the theorem-like environments the retained text needs and adds the article's own macro definitions that the retained text needs (\texttt{\textbackslash ISO}, \texttt{\textbackslash One}, \texttt{\textbackslash R}, \texttt{\textbackslash bZ}), with extra wrapper packages (bm, bbm, dsfont, xspace, cancel, mathtools). The delivered numbering is the unit's own. Assets: no retained span contains a figure environment, a graphics-inclusion command, a PDF-page inclusion, a table or a TikZ picture; no image file is copied into this package, the package declares no asset dependency and no image file is redistributed with it. Licence: version arXiv:2501.06133v3 of this article is distributed under the Creative Commons Attribution 4.0 International licence, as recorded in the arXiv licence metadata for this exact version and shown on the abstract page https://arxiv.org/abs/2501.06133v3. A full rights notice is delivered at licenses/arxiv-2501.06133-CC-BY-4.0.txt. Characters: every retained excerpt is pure ASCII, so the delivered engine, XeLaTeX with its default Latin Modern text font, reports no missing character.
 \begin{equation}\label{Eq:utildeiso}
     \tilde{u}_i = \sum_{t=1}^{m-1}  r_i\,\One{n_t\leq i<n_{t+1}},
 \end{equation}

\begin{lemma}\label{lem:a-median-result}
Suppose $a\in\R^m$ and $b\in\R^n$. Let $c=(a,b)\in\R^{m+n}$ be the concatenation of the two vectors. If $\textnormal{Med}(c)< \textnormal{Med}(a)$, then $\textnormal{Med}(c)\geq\textnormal{Med}(b)$; similarly, if $\textnormal{Med}(c)>\textnormal{Med}(a)$, then $\textnormal{Med}(c)\leq\textnormal{Med}(b)$.
\end{lemma}

\begin{lemma}\label{lem:characterizing-isotonic-L1-projection_}
In the representation~\eqref{Eq:utildeiso} of the vector $\tilde{u}$, for each $t\in[m-1]$,     \begin{equation}\label{eq:cumulative-median-equals}
    r_t = \textnormal{Med}\bigl((u_\ell)_{n_t \leq \ell \leq n_{t+1}-1}\bigr) ,
    \end{equation}
    and moreover, for each $i \in \{n_t,n_t+1,\ldots,n_{t+1}-1\}$,
\begin{equation}\label{eq:cumulative-median-inequalities}
    \textnormal{Med}\bigl((u_\ell)_{\ell:Z_{n_t}\leq Z_\ell\leq Z_i}\bigr) 
        \geq r_t \geq  \textnormal{Med}\bigl((u_\ell)_{\ell:Z_i\leq Z_\ell\leq Z_{n_{t+1}-1}}\bigr).    
    \end{equation}
\end{lemma}

\begin{proof}
    First, since $\tilde{u}\in\widehat{\mathcal{C}}_\ISO(\bZ)$, and $\tilde{u}_{n_t-1}<\tilde{u}_{n_t}$ for each $t=2,\dots,m-1$, we have 
    $Z_{n_t-1} < Z_{n_t}$ for each $t=2,\dots,m-1$ and consequently
    \[n_t \leq \ell\leq n_{t+1}-1 \ \Longleftrightarrow \ Z_{n_t}\leq Z_\ell \leq Z_{n_{t+1}-1}\]
    for each $t\in[m-1]$.
    Therefore, establishing~\eqref{eq:cumulative-median-equals} is equivalent to proving that
    \[r_t = \textnormal{Med}\bigl((u_\ell)_{\ell:Z_{n_t} \leq Z_\ell \leq Z_{n_{t+1}-1}}\bigr),\]
    which is an immediate consequence of the inequalities in~\eqref{eq:cumulative-median-inequalities} (applied with $i=n_{t+1}-1$ on the left and with $i=n_t$ on the right). Therefore we now only need to prove these inequalities.

    \paragraph{Upper bound on $r_t$.}
    First we will prove that
    \begin{equation}\label{eq:cumulative-median-inequalities-1}
    r_t\leq  \textnormal{Med}\bigl((u_\ell)_{\ell:Z_{n_t}\leq Z_\ell\leq Z_i}\bigr)
\end{equation}
for all $t\in [m-1]$ and all $i\in\{n_t,\dots,n_{t+1}-1\}$.
By definition, we have 
\[r_t = \tilde{u}_{n_t} = \max_{j\in [n]: Z_j\leq Z_{n_t}}\min_{k\in [n]: Z_k\geq Z_{n_t}} \textnormal{Med}\bigl((u_\ell)_{\ell:Z_j\leq Z_\ell\leq Z_k}\bigr).\]
As noted above, $Z_j < Z_{n_t}$ for any $j<n_t$. 
Therefore we can write
\[
\{j\in[n]: Z_j \leq Z_{n_t}\} = \{j\in[n]: j< n_t\} \cup \{j\in[n]: Z_j = Z_{n_t}\},
\]
and hence,
\[
r_t = \max\Big\{\underbrace{\max_{j < n_t}\min_{k\in [n]: Z_k\geq Z_{n_t}} \textnormal{Med}\bigl((u_\ell)_{\ell:Z_j\leq Z_\ell\leq Z_k}\bigr)}_{=:r_{t,1}} , \ \underbrace{\min_{k\in [n]: Z_k\geq Z_{n_t}} \textnormal{Med}\bigl((u_\ell)_{\ell:Z_{n_t}\leq Z_\ell\leq Z_k}\bigr)}_{=:r_{t,2}}\Big\},
\]
where we take the convention that the maximum over an empty set is equal to $-\infty$ (to handle the case $t=1$, i.e., $\{j\in[n]:j<n_1\}=\emptyset$, so that $r_{t,1}=-\infty$ and $r_t = r_{t,2}$. 
If instead $t\geq 2$, then we claim that $r_{t,1}\leq r_{t,2}$, so that $r_t = r_{t,2}$ again. Suppose for a contradiction that $r_{t,1}>r_{t,2}$. Let $\hat{j}\in [n_{t}-1]$ and $\hat{k}\in [n]$ be such that $Z_{\hat{k}} \geq Z_{n_t}$ and
\begin{equation}\label{eq:expression-rt1-rt2}
    r_{t,1}=\min_{k\in [n]: Z_k\geq Z_{n_t}} \textnormal{Med}\bigl((u_\ell)_{\ell:Z_{\hat{j}}\leq Z_\ell\leq Z_k}\bigr)\, \text{ and }\,r_{t,2}=\textnormal{Med}\bigl((u_\ell)_{\ell:Z_{n_t}\leq Z_\ell\leq Z_{\hat{k}}}\bigr).
\end{equation}
Then, $\textnormal{Med}\bigl((u_\ell)_{\ell:Z_{\hat{j}}\leq Z_\ell\leq Z_{\hat{k}}}\bigr)\geq r_{t,1}>\textnormal{Med}\bigl((u_\ell)_{\ell:Z_{n_t}\leq Z_\ell\leq Z_{\hat{k}}}\bigr)$.
Since $\{j: j<n_t\} = \{j: Z_j\leq Z_{n_t-1}\}$ (due to $Z$ being a nondecreasing vector with $Z_{n_t-1}<Z_{n_t}$ as established above), it holds that $Z_{\hat{j}} \leq Z_{n_t-1} <Z_{n_t}$ and thus, by Lemma~\ref{lem:a-median-result},
\[
\textnormal{Med}\bigl((u_\ell)_{\ell:Z_{\hat{j}}\leq Z_\ell\leq Z_{\hat{k}}}\bigr)\leq \textnormal{Med}\bigl((u_\ell)_{\ell:Z_{\hat{j}}\leq Z_\ell\leq Z_{n_t-1}}\bigr).
\]
Now, we also have
\begin{equation}\label{eq:bounding-rt1-below}
    r_{t,1}=r_t>\tilde{u}_{n_t-1}\geq \min_{k\in [n]: Z_k\geq Z_{n_t-1}} \textnormal{Med}\bigl((u_\ell)_{\ell:Z_{\hat{j}}\leq Z_\ell\leq Z_k}\bigr)
\end{equation}
Moreover, by noting that $\{k : Z_k \geq Z_{n_t-1}\}=\{k : Z_k \geq Z_{n_t}\}\cup \{k:Z_k=Z_{n_t-1}\}$ and by~\eqref{eq:expression-rt1-rt2}, it follows that 
\begin{align*}
    &\min_{k\in [n]: Z_k\geq Z_{n_t-1}} \textnormal{Med}\bigl((u_\ell)_{\ell:Z_{\hat{j}}\leq Z_\ell\leq Z_k}\bigr)\\&\hspace{1.6cm}=\min\Bigl\{\min_{k\in [n]: Z_k\geq Z_{n_t}} \textnormal{Med}\bigl((u_\ell)_{\ell:Z_{\hat{j}}\leq Z_\ell\leq Z_k}\bigr), \textnormal{Med}\bigl((u_\ell)_{\ell:Z_{\hat{j}}\leq Z_\ell\leq Z_{n_t-1}}\bigr)\Bigr\}\\
    &\hspace{1.6cm}=\min\Bigl\{r_{t,1}, \textnormal{Med}\bigl((u_\ell)_{\ell:Z_{\hat{j}}\leq Z_\ell\leq Z_{n_t-1}}\bigr)\Bigr\}= \textnormal{Med}\bigl((u_\ell)_{\ell:Z_{\hat{j}}\leq Z_\ell\leq Z_{n_t-1}}\bigr),
\end{align*}
where the last step follows by~\eqref{eq:bounding-rt1-below}. Thus, we have that
\[\textnormal{Med}\bigl((u_\ell)_{\ell:Z_{\hat{j}}\leq Z_\ell\leq Z_{\hat{k}}}\bigr)\geq r_{t,1}> \textnormal{Med}\bigl((u_\ell)_{\ell:Z_{\hat{j}}\leq Z_\ell\leq Z_{n_t-1}}\bigr),\]
which is a contradiction.
Thus we again have $r_t=r_{t,2}$ in this case. Therefore,
for all $t\in [m-1]$,
\[r_t = r_{t,2} = \min_{k\in [n]: Z_k\geq Z_{n_t}} \textnormal{Med}\bigl((u_\ell)_{\ell:Z_{n_t}\leq Z_\ell\leq Z_k}\bigr) \leq \min_{i\in\{n_t,\dots,n_{t+1}-1\}}\textnormal{Med}\bigl((u_\ell)_{\ell:Z_{n_t}\leq Z_\ell\leq Z_i}\bigr).\]
This establishes the desired upper bound~\eqref{eq:cumulative-median-inequalities-1} on $r_t$.

    \paragraph{Lower bound on $r_t$.}
    Next we will prove that
    \begin{equation}\label{eq:cumulative-median-inequalities-2}
    r_t\geq  \textnormal{Med}\bigl((u_\ell)_{\ell:Z_i\leq Z_\ell\leq Z_{n_{t+1}-1}}\bigr)
\end{equation}
for all $t\in[m-1]$ and all $i\in\{n_t,\dots,n_{t+1}-1\}$. This proof is not simply a symmetric version of the proof of the upper bound, since the vector $\tilde{u}$ is defined by taking the maximum of the minimum (and not vice versa).

By definition, we have 
\[r_t = \tilde{u}_{n_{t+1}-1} = \max_{j\in [n]: Z_j\leq Z_{n_{t+1}-1}}\min_{k\in [n]: Z_k\geq Z_{n_{t+1}-1}} \textnormal{Med}\bigl((u_\ell)_{\ell:Z_j\leq Z_\ell\leq Z_k}\bigr).\]
If $t=m-1$ then $n_{t+1}-1=n$, and so $\{k\in[n] : Z_k \geq Z_{n_{t+1}-1}\} = \{k\in[n]: Z_k \geq Z_n\}$. Therefore, for $t=m-1$ we have
\begin{align*}
r_t = \max_{j\in [n]: Z_j\leq Z_n}\textnormal{Med}\bigl((u_\ell)_{\ell:Z_j\leq Z_\ell\leq Z_n}\bigr) &= \max_{j \in [n]} \textnormal{Med}\bigl((u_\ell)_{\ell:Z_j\leq Z_\ell\leq Z_n}\bigr) \\
&\geq \max_{j \in\{n_t,\dots,n\}}\textnormal{Med}\bigl((u_\ell)_{\ell:Z_j\leq Z_\ell\leq Z_{n}}\bigr).
\end{align*}

This proves the claim for the case $t=m-1$.

From this point on we assume $t\in[m-2]$. 
As noted above, $Z_k > Z_{n_{t+1}-1}$ for any $k\geq n_{t+1}$. 
Therefore we can write
\[
\{k\in[n]: Z_k \geq Z_{n_{t+1}-1}\} = \{k\in[n]: k\geq n_{t+1}\} \cup \{k\in[n]: Z_k = Z_{n_{t+1}-1}\},
\]
and hence,
\[
r_t = \max_{j\in[n]:Z_j \leq Z_{n_{t+1}-1}}\min\Big\{\underbrace{\min_{k\geq n_{t+1}}\textnormal{Med}\bigl((u_\ell)_{\ell:Z_j\leq Z_\ell\leq Z_k}\bigr)}_{=:r_{t,j,1}} , \ \underbrace{ \textnormal{Med}\bigl((u_\ell)_{\ell:Z_j\leq Z_\ell\leq Z_{n_{t+1}-1}}\bigr)}_{=:r_{t,j,2}}\Big\}.
\]
Fix any $j$ with $Z_j \leq Z_{n_{t+1}-1}$, and suppose that $r_{t,j,1} < r_{t,j,2}$. Then
\[
\textnormal{Med}\bigl((u_\ell)_{\ell:Z_j\leq Z_\ell\leq Z_k}\bigr) < \textnormal{Med}\bigl((u_\ell)_{\ell:Z_j\leq Z_\ell\leq Z_{n_{t+1}-1}}\bigr)
\]
for some $k\geq n_{t+1}$. Then by Lemma~\ref{lem:a-median-result}, we must have
 \[\textnormal{Med}\bigl((u_\ell)_{\ell:Z_j\leq Z_\ell\leq Z_k}\bigr) \geq \textnormal{Med}\bigl((u_\ell)_{\ell:Z_{n_{t+1}-1}< Z_\ell\leq Z_k}\bigr).\]
Now let $t' \geq t+1$ be such that $n_{t'}\leq k < n_{t'+1}$.  Since $Z_{n_{t+1}-1} < Z_{n_{t+1}}$ as explained above, we have by Lemma~\ref{lem:a-median-result} that 
\begin{multline*}
\textnormal{Med}\bigl((u_\ell)_{\ell:Z_{n_{t+1}-1}< Z_\ell\leq Z_k}\bigr) = \textnormal{Med}\bigl((u_\ell)_{\ell:Z_{n_{t+1}}\leq Z_\ell\leq Z_k}\bigr) \\\geq \min\left\{\textnormal{Med}\bigl((u_\ell)_{\ell:Z_{n_{t+1}}\leq Z_\ell < Z_{n_{t+2}}}\bigr), \dots,\textnormal{Med}\bigl((u_\ell)_{\ell:Z_{n_{t'-1}}\leq Z_\ell < Z_{n_{t'}}}\bigr),\textnormal{Med}\bigl((u_\ell)_{\ell:Z_{n_{t'}}\leq Z_\ell \leq Z_k}\bigr)\right\}\\\geq \min\{r_{t+1},\dots,r_{m-1}\} = r_{t+1},
\end{multline*}
where the second inequality applies the upper bound~\eqref{eq:cumulative-median-inequalities-1} on each $r_{t+1},\dots,r_{m-1}$. But we also have
\[
r_t \geq \min\{r_{t,j,1},r_{t,j,2}\} = r_{t,j,1} \geq r_{t+1},
\]
which is a contradiction.

Therefore we have $r_{t,j,1}\geq r_{t,j,2}$ for all $j$ such that $Z_j \leq Z_{n_{t+1}-1}$, and so
\begin{multline*}
r_t = \max_{j\in[n]:Z_j \leq Z_{n_{t+1}-1}}r_{t,j,2} = \max_{j\in[n]:Z_j \leq Z_{n_{t+1}-1}} \textnormal{Med}\bigl((u_\ell)_{\ell:Z_j\leq Z_\ell\leq Z_{n_{t+1}-1}}\bigr) \\\geq \max_{j\in\{n_t,\dots,n_{t+1}-1\}}\textnormal{Med}\bigl((u_\ell)_{\ell:Z_j\leq Z_\ell\leq Z_{n_{t+1}-1}}\bigr).
\end{multline*}
This proves the lower bound~\eqref{eq:cumulative-median-inequalities-2}, and thus completes the proof of the lemma.
\end{proof}

\end{document}
